\section{The story begins: what makes navigation social?}
\label{sec:fundamentals}

\subsection{From reaching a coordinate to sharing a space}
The simplest navigation question is geometric: can a robot move from its present
configuration to a goal without intersecting an obstacle? If the map and robot
shape are fixed, a collision-free path is a sequence of configurations in free
space. Even this question is not the whole control problem. A path does not specify
how fast to traverse it, whether a motor can produce the requested acceleration,
or whether the robot can stop before a previously unseen obstacle appears.
DWA made this distinction explicit long before modern deep learning: searching
reachable, stoppable velocities is different from selecting an appealing direction
and converting it into motion afterward \citep{dwa}.

People introduce a second distinction. A person has an intention, a perception of
the robot and the ability to react. A route that is physically possible may force
someone to change direction, interrupt a group, block access or produce an
uncomfortable encounter. Conversely, the most conservative geometric solution may
stop indefinitely in the middle of the shared space. Thus neither ``reaches the
goal'' nor ``never moves into an occupied cell'' is a complete definition of
successful social navigation. The surveys organize this problem around human
awareness, interaction context and the purpose of the task
\citep{survey2013,survey2024}.

The distinction is especially important for a delivery humanoid. Independent
passage through a corridor, approaching a recipient and carrying an object past
a seated person are different social situations. A distance acceptable during an
invited handover may be inappropriate when overtaking an unaware pedestrian.
The desired behavior therefore depends on the robot's task and the human's role,
not only on relative coordinates. This is why the first research decision must
be a task definition, rather than a choice between an optimization-based and a
learning-based planner.

\subsection{Five objectives that should not be collapsed}
For the present review, the navigation problem has five distinct outcome families.
\begin{enumerate}
\item \textbf{Physical integrity:} avoid prohibited contacts, loss of balance,
falls and dropped payloads. The relevant geometry includes whatever can hit a
person, not necessarily just the base footprint.
\item \textbf{Task completion:} reach the intended goal and satisfy its stopping
or delivery condition. Briefly entering a goal region is not necessarily a
completed docking or placement action.
\item \textbf{Liveness and efficiency:} continue making useful progress, resolve
temporary blockage and complete the task without excessive delay or intervention.
\item \textbf{Social appropriateness:} produce behavior that is understandable,
contextually suitable and acceptable to affected people.
\item \textbf{Evidence validity:} establish that the sensors, scorer, timestamps
and event definitions support the claimed outcome.
\end{enumerate}
These outcomes can conflict. Increased clearance may lengthen a route; abrupt
avoidance may protect base clearance while destabilizing a biped; early stopping
may reduce the robot's motion into a person but increase the time it obstructs a
passage. The evaluation guidelines consequently advocate a multidimensional
account rather than one unqualified social score \citep{guidelines}.

It is useful to distinguish a \emph{proxy} from its target construct. Time within
a specified personal-space radius is a geometric statistic. Perceived discomfort
is a human response. The statistic may correlate with discomfort under a declared
condition, but the relationship has to be tested. The proxemics literature makes
the dependence on human, robot, environment and context particularly explicit
\citep{proxemics2026}. The same caution applies to smoothness: low command jerk
does not prove that the executed motion is smooth, and smooth motion does not
prove that its social intent is legible.

\subsection{A decision model that exposes the hidden assumptions}
Let $x_t$ denote the robot's physical state, $z_t$ the surrounding agents and
environment, and $o_{0:t}$ the available observation history. A planner chooses
an action using a belief or estimate rather than complete knowledge:
\begin{equation}
u_t=\pi(b_t,g),\qquad b_t=\mathcal{E}(o_{0:t}),\qquad
x_{t+1}=f(x_t,u_t,z_t,w_t),
\end{equation}
where $g$ is the task goal and $w_t$ represents disturbances and modeling error.
This notation is explanatory, not a new algorithm. It identifies four questions
that recur throughout the literature: what is observed, how uncertainty is
represented, what action the policy actually commands, and how that action changes
the physical state.

A planar simulator may make $u_t$ the next velocity and realize it immediately.
A humanoid navigation stack may instead make $u_t$ a desired body velocity,
then pass it through normalization, acceleration limits, a locomotion network,
joint targets, PD actuation and contact dynamics. The physical transition is
therefore mediated by a second policy and by execution state. Comparing these
systems only on the label ``velocity action'' hides a major difference.
SARL explicitly uses instantaneous velocity achievement; biped MPC explicitly
models balance and terminal stopping; a LiDAR-to-motor policy integrates another
part of the chain \citep{sarl,biped2016,safehumanoid}.

People's behavior introduces an additional conditional dependence:
\begin{equation}
p(z_{t+1:t+H}\mid o_{0:t},u_{t:t+H-1}).
\end{equation}
An independent constant-velocity forecast omits the dependence on robot action.
Reciprocal avoidance imposes a particular response rule. An interaction-conditioned
predictor tries to learn that dependence. None is appropriate merely because it
is more complicated; the correct choice depends on whether the surrounding
people actually observe and adapt to the robot and on the information available
to infer their response \citep{orca,trajectron,prediction}.

\subsection{What geometric methods solved, and what they left open}
Velocity-space methods turned dynamic avoidance into a tractable geometric
problem. DWA restricts the feasible command set using acceleration and stopping;
velocity obstacles reason about relative velocities leading to future collision;
ORCA distributes corrective responsibility among cooperating agents
\citep{dwa,vo,orca}. Timed trajectory optimization adds an explicit temporal
structure to the local path \citep{teb}. These are not obsolete baselines simply
because later methods learn their representations. They expose assumptions and
failure reasons that are often harder to inspect in a neural policy.

Their limitations are specific rather than universal. A reciprocal guarantee
does not apply unchanged to a person who never reacts. A safe velocity can become
unsafe after post-hoc clamping. A finite-horizon collision-free trajectory may
have no safe continuation. A local cost optimum may not escape a blocked route.
The ORCA chapter itself recognizes the loss of its strict guarantee after
kinematic/dynamic post-processing \citep{orca}. A serious comparison should use
these limitations as experimental variables, not as a rhetorical reason to omit
classical methods.

\subsection{Why cooperation and freezing are two sides of the same problem}
Suppose a robot predicts several plausible trajectories for every nearby person
and refuses to enter any of them. As crowd density and prediction horizon grow,
the union of those regions can cover every useful route. The robot freezes, even
though people could pass each other by making small coordinated adaptations.
Models of cooperation address this by reasoning about mutually compatible
behavior \citep{freezing}. Social-force and reciprocal models likewise generate
collective organization from local interactions, although they encode different
behavioral assumptions \citep{socialforce,orca}.

The opposite failure is overconfidence in cooperation. A controller may select
a path that is safe only if a pedestrian slows or steps aside. An inattentive
person, mobility-aid user or person approaching from behind may not do so. The
interesting research problem is therefore not to choose universally between
cooperation and conservatism. It is to define which assumptions justify a
particular action and what the robot does when the assumptions are uncertain or
contradicted by observations. Nonreactive actors are useful for testing unilateral
avoidance, but they cannot evaluate the quality of negotiated interaction.

\subsection{Why learning changed the representation of the crowd}
Learned navigation initially made it possible to amortize expensive interaction
reasoning into a value function. CADRL represents local interactions in a
robot-centric state; SARL learns how to aggregate several people; RGL learns
relations and future state evolution; DS-RNN adds structured temporal reasoning
\citep{cadrl,sarl,rgl,dsrnn}. This progression is more informative than a simple
``classical versus AI'' division. The models move different parts of inference
and planning into learned components and make different assumptions about the
state available at deployment.

Learning also changes the experimental burden. A method's behavior depends on
the training distribution, crowd policy, reward, privileged inputs and model
selection procedure. A planner trained with circle-crossing ORCA agents can
learn useful anticipation while remaining poorly characterized under detector
dropout or a different walking controller. Conversely, a failure after naive
transfer may measure an interface mismatch rather than an inherent weakness of
the learning architecture. Equal input information and comparable tuning budgets
are prerequisites for an interpretable comparison.

\subsection{Prediction is an input to a decision, not the final objective}
Social LSTM, Social GAN and Trajectron++ demonstrate increasingly expressive
representations of pedestrian futures: interaction pooling, multimodal sampling,
heterogeneous data and dynamics/ego-plan conditioning
\citep{sociallstm,socialgan,trajectron}. Yet a robot must decide before knowing
which future will occur. A predictor's best sample can be close to the recorded
trajectory while the distribution assigns too little probability to a dangerous
alternative. Conversely, a simple short-horizon estimate can support an effective
decision when observations are frequent and the robot has time to brake.

This motivates a decision-centered evaluation: which errors change the chosen
action, whether the predicted occupied region covers the actual person at the
encounter, and whether the resulting robot trajectory is executable. Lower
average displacement error is useful evidence about a forecasting task, but is
not by itself evidence of fewer contacts or better social behavior. The later
sections use this distinction to connect perception, timing and control.

\subsection{Embodiment closes the loop}
A biped must remain balanced while it avoids people. The 2016 biped crowd paper
already unifies balance, stopping and perception uncertainty through viability
and lexicographic MPC \citep{biped2016}. Learned locomotion later broadened the
achievable repertoire through adaptation, perceptive control and capability-aware
navigation \citep{rma,rough,parkour,wheeledleg}. Humanoid imitation methods show
how arm objectives, posture and dynamics mismatch alter the control problem
\citep{expressive,asap}. These developments explain why a planar command model
cannot simply be assumed correct for a humanoid.

They also constrain novelty. ``Account for dynamics'' and ``combine navigation
with locomotion'' are already established research directions. What can remain
open is a particular, validated interface: for example, how a measured response
distribution for a deployed policy should combine with observation age and
person uncertainty when selecting an avoidance action. The contribution must
be the specified model, decision rule or causal result, rather than the fact that
the robot has legs.

\subsection{From a successful maneuver to dependable social autonomy}
Recent LiDAR-conditioned policies and visual generative navigation systems
already demonstrate sensor-driven avoidance and transfer to legged or humanoid
platforms \citep{safehumanoid,omni,flux}. Benchmarks now include grounded human
trajectories, interactive avatars, diverse urban scenes and composed tasks
\citep{socnavbench,habitat3,metaurban,omninav,dped}. A research programme must
therefore ask why successful components still fail when composed, rather than
assuming that no one has integrated them before.

The central story of this review is that \textbf{the field progressively exposes
more of the closed loop}. Geometry makes collision conditions explicit;
interaction models make human response explicit; learning makes richer behavior
representations practical; perception makes partial observability unavoidable;
embodied control makes command execution conditional; and complete-mission
evaluation exposes failure propagation across all of them. The proposed novelty
roadmap follows that story: establish a measured baseline, define a testable new
interface or mechanism, isolate its effect and demonstrate that it matters under
conditions the nearest prior methods do not already resolve.
